\documentclass[a4paper,11pt]{article}
\usepackage[T2A]{fontenc}
\usepackage[utf8]{inputenc}
\usepackage[russian]{babel}
\usepackage{amsmath,amssymb,mathtools}
\usepackage{geometry}
\usepackage{helvet}
\renewcommand{\familydefault}{\sfdefault}
\usepackage{microtype}
\usepackage{array,booktabs,tabularx}
\usepackage{xcolor}
\usepackage{hyperref}
\usepackage{tikz}
\usepackage{fancyhdr}
\usepackage{enumitem}
\geometry{left=20mm,right=19mm,top=18mm,bottom=19mm}
\definecolor{accent}{HTML}{006B6C}
\definecolor{ink}{HTML}{242C32}
\definecolor{muted}{HTML}{56616C}
\definecolor{error}{HTML}{A42B36}
\hypersetup{colorlinks=true,linkcolor=accent,urlcolor=accent,pdfauthor={Лёвин Артём Александрович},pdftitle={Ревью домашней работы — Володя Хачатурян, 15.11.2026}}
\setlength{\parindent}{0pt}
\setlength{\parskip}{5pt}
\setlength{\tabcolsep}{4pt}
\renewcommand{\arraystretch}{1.18}
\setlist[itemize]{topsep=2pt,itemsep=1pt,leftmargin=*}
\pagestyle{fancy}
\fancyhf{}
\fancyhead[L]{\footnotesize\color{muted} Ревью домашней работы}
\fancyhead[R]{\footnotesize\color{muted} Движение по окружности}
\fancyfoot[C]{\footnotesize\color{muted}\thepage}
\renewcommand{\headrulewidth}{0pt}
\renewcommand{\footrulewidth}{0pt}
\newcommand{\keyidea}{\textcolor{accent}{\textbf{Ключевая идея.}}}
\newcommand{\mistake}{\textcolor{error}{\textbf{Ошибка.}}}
\newcommand{\result}{\textcolor{accent}{\textbf{Итог.}}}
\newcommand{\taskhead}[1]{\phantomsection\label{task:#1}\begin{center}{\Large\bfseries\color{ink}Задание №\,#1}\end{center}\vspace{3pt}}
\newcommand{\field}[1]{\par\vspace{2pt}\textbf{#1}\quad}
\begin{document}
\thispagestyle{empty}
\vspace*{36mm}
\begin{center}
{\large\bfseries\color{accent} ФИЗИКА}\\[18mm]
{\Huge\bfseries\color{ink} Ревью домашней работы}\\[9mm]
{\LARGE Движение по окружности}\\[20mm]
{\large Володя Хачатурян}\\[14mm]
{\normalsize 15 ноября 2026 г.}\\[32mm]
{\normalsize Лёвин Артём Александрович}\\[3mm]
{\small\color{muted} эксклюзивно для Володи Хачатуряна}
\end{center}
\vfill
\begin{center}
\begin{tikzpicture}
\draw[accent!50,line width=0.8pt] (0,0) .. controls (2.4,0.6) and (4.4,-0.6) .. (7,0)
  .. controls (9.6,0.55) and (11.2,-0.35) .. (13,0.1);
\end{tikzpicture}
\end{center}
\clearpage
\begin{center}{\Large\bfseries\color{ink}Сводная таблица}\end{center}
\noindent\textcolor{muted}{\small Проверены только задания, для которых на фотографиях видны записи. Верно решённые задачи в таблицу не включаются.}\par
\vspace{4mm}
{\footnotesize
\begin{tabularx}{\linewidth}{@{}>{\centering\arraybackslash}p{9mm}>{\raggedright\arraybackslash}p{20mm}>{\raggedright\arraybackslash}X>{\raggedright\arraybackslash}p{20mm}>{\raggedright\arraybackslash}p{20mm}@{}}
\toprule
\textbf{№} & \textbf{Тема} & \textbf{Суть ошибки} & \textbf{Ваш ответ} & \textbf{Верный ответ}\\
\midrule
\hyperref[task:4]{4} & Вершина моста & Данные другой задачи. Не учтена радиальная сила. & 40 Н, 12 Н & 3600 Н\\
\hyperref[task:5]{5} & Потеря контакта & Рассмотрена сила трения вместо скорости отрыва. & Масса, кг & 30 м/с\\
\hyperref[task:6]{6} & Ведёрко вверху & Иные исходные данные, натяжение не вычислено. & Ускорение & 16 Н\\
\hyperref[task:9]{9} & Ведёрко внизу & Уравнение и итоговый ответ отсутствуют. & Нет ответа & 12 м/с\\
\hyperref[task:10]{10} & Прочность нити & Нет уравнений и сравнения с предельной силой. & Нет ответа & Порвётся внизу\\
\bottomrule
\end{tabularx}
}\par
\vspace{5mm}
\textbf{Замечание к исходным материалам.} Фотографии \texttt{01.jpg} и \texttt{02.jpg} совпадают; \texttt{03.jpg} и \texttt{04.jpg} также совпадают. Записи №\,4--6 заметно расходятся с печатным листом условий. Задания №\,1--3 и №\,7--8 на представленных страницах отсутствуют; их выполнение оценить нельзя.

\textbf{Принято при расчётах:} $g=10\,\text{м/с}^2$. Значение ускорения свободного падения на печатном листе отдельно не задано.
\clearpage
\thispagestyle{empty}
\vspace*{0.36\textheight}
\begin{center}{\Huge\bfseries\color{accent}Разбор ошибок}\end{center}
\clearpage

\taskhead{4}
\field{Условие.} В верхней точке выпуклого участка дороги автомобиль массой $600$ кг движется со скоростью $10$ м/с. Радиус кривизны $25$ м. Ось направлена вниз. Найти реакцию опоры $N$.

\field{Ваше решение.} В записи используются $N=mg=4\cdot10=40$ Н и формула силы трения $F_{\text{тр}}=\mu N=0{,}3\cdot40=12$ Н. Эти данные и коэффициент трения отсутствуют в напечатанном условии №\,4.

\field{Ваш ответ.} $N=40$ Н; дополнительно $F_{\text{тр}}=12$ Н.

\field{Где именно возникает ошибка.} Последний пригодный элемент --- идея использовать силу тяжести $mg$. Первый неверный для данной задачи переход: $N=mg$. При движении по криволинейной траектории результирующая радиальная сила отлична от нуля. Подстановка массы $4$ кг также не соответствует условию ($600$ кг).

\mistake{} Реакция опоры ошибочно приравнена весу покоящегося тела; дополнительно использованы посторонние исходные данные.

\field{Почему это неверно.} В верхней точке выпуклого моста центр кривизны расположен \textbf{снизу}. Сила тяжести $mg$ направлена вниз, реакция опоры $N$ --- вверх. При выбранной положительной оси вниз их проекции имеют противоположные знаки. Следовательно, результирующая сила $mg-N$ создаёт центростремительное ускорение $v^2/R$.

\keyidea{} В верхней точке выпуклого моста $mg-N=mv^2/R$; поэтому при $v>0$ реакция опоры меньше $mg$.

\field{Исправление.} Сохраните радиальное направление оси вниз и запишите второй закон Ньютона в проекциях: $mg-N=mv^2/R$.

\field{Полное правильное решение.}
\begin{align*}
mg-N&=m\frac{v^2}{R},\\
N&=mg-m\frac{v^2}{R},\\
N&=600\cdot10-600\cdot\frac{10^2}{25}\\
 &=6000-2400=\boxed{3600\ \text{Н}}.
\end{align*}
\field{Проверка.} $0<N<mg$: $0<3600<6000$ Н. Кроме того, $6000-3600=2400$ Н и $mv^2/R=2400$ Н.

\result{} Реакция опоры $N=3600$ Н.

\field{Задача на закрепление.} Автомобиль массой $800$ кг проходит вершину выпуклого моста со скоростью $10$ м/с. Радиус кривизны $40$ м. Найдите реакцию опоры, приняв $g=10\,\text{м/с}^2$.
\clearpage

\taskhead{5}
\field{Условие.} Найти наибольшую скорость автомобиля в верхней точке выпуклого моста радиусом $90$ м, при которой контакт с дорогой ещё сохраняется.

\field{Ваше решение.} В тетради вместо сил в верхней точке моста рассматриваются сила трения, коэффициент трения $\mu=0{,}2$ и вычисление массы. Найти скорость в метрах в секунду по этим записям не удалось.

\field{Ваш ответ.} Значение массы в килограммах; ответ на вопрос о скорости отсутствует.

\field{Где именно возникает ошибка.} Первый неверный шаг --- выбор модели с трением $F_{\text{тр}}=\mu N$. Печатное условие требует найти границу контакта на выпуклом участке, поэтому определяющим является обращение реакции опоры в нуль.

\mistake{} Выбрана сила, не определяющая предельную скорость в этой постановке, и потеряна искомая физическая величина.

\field{Почему это неверно.} Пока автомобиль соприкасается с дорогой, опора может толкать его вверх с силой $N\ge0$. При росте скорости требуемая направленная вниз результирующая сила возрастает, поэтому $N$ уменьшается. На границе отрыва $N=0$. Отрицательная реакция опоры означала бы, что дорога должна притягивать автомобиль; обычная опора этого делать не может.

\keyidea{} Для предельной скорости на вершине выпуклого моста одновременно выполняются $mg-N=mv^2/R$ и $N=0$.

\field{Исправление.} Запишите радиальное уравнение и положите $N=0$.

\field{Полное правильное решение.}
\begin{align*}
mg-N&=m\frac{v^2}{R},\qquad N=0,\\
mg&=m\frac{v^2}{R},\\
g&=\frac{v^2}{R},\\
v_{\max}&=\sqrt{gR}=\sqrt{10\cdot90}
=\sqrt{900}=\boxed{30\ \text{м/с}}.
\end{align*}
\field{Проверка.} При $v=30$ м/с величина $mv^2/R$ равна $mg$, значит $N=0$. При меньших скоростях $N>0$, при больших формула потребовала бы $N<0$.

\result{} Максимальная скорость без потери контакта --- $30$ м/с.

\field{Задача на закрепление.} Радиус выпуклого моста $40$ м. Найдите максимальную скорость автомобиля на его вершине без отрыва от дороги. Примите $g=10\,\text{м/с}^2$.
\clearpage

\taskhead{6}
\field{Условие.} Ведёрко массой $2$ кг равномерно вращается на нити в вертикальной плоскости. В верхней точке радиус $R=2$ м, скорость $v=6$ м/с. Найти натяжение нити.

\field{Ваше решение.} В тетради указаны иные данные, в том числе масса $5$ кг, коэффициент трения $0{,}25$ и сила $20$ Н; вычисляется ускорение через разность сил. В исходной задаче трения нет, требуется натяжение нити.

\field{Ваш ответ.} Получено ускорение для другого условия; искомое натяжение не указано.

\field{Где именно возникает ошибка.} До расчётов не определено направление к центру окружности и не записан баланс сил для верхней точки. Выбор схемы с силой трения не позволяет решить данную задачу.

\mistake{} В уравнении движения не отражены две направленные вниз силы: тяжесть и натяжение нити.

\field{Почему это неверно.} В верхней точке центр окружности находится внизу. Земля притягивает ведёрко вниз, и натянутая нить тоже тянет его к центру --- вниз. Поэтому радиальные проекции этих сил \textbf{складываются}. Нельзя использовать формулу для нижней точки, где тяжесть направлена от центра.

\keyidea{} В верхней точке вращения на нити $T+mg=mv^2/R$, следовательно, $T=mv^2/R-mg$.

\field{Исправление.} Выберите ось вниз и включите в неё обе силы с положительными знаками.

\field{Полное правильное решение.}
\begin{align*}
T+mg&=m\frac{v^2}{R},\\
T&=m\frac{v^2}{R}-mg,\\
T&=2\cdot\frac{6^2}{2}-2\cdot10\\
 &=36-20=\boxed{16\ \text{Н}}.
\end{align*}
\field{Проверка.} $T=16$ Н положительно: нить может оставаться натянутой. Сумма направленных к центру сил равна $16+20=36$ Н, что совпадает с $mv^2/R$.

\result{} Натяжение нити в верхней точке $16$ Н.

\field{Задача на закрепление.} Ведёрко массой $1{,}5$ кг вращается на нити длиной $3$ м. В верхней точке его скорость $8$ м/с. Найдите натяжение нити при $g=10\,\text{м/с}^2$.
\clearpage

\taskhead{9}
\field{Условие.} В нижней точке вертикальной окружности ведёрко массой $2$ кг движется на нити длиной $4$ м. Максимально допустимое натяжение --- $92$ Н. Найти наибольшую допустимую скорость.

\field{Ваше решение.} На листе под номером 9 приведены отдельные записи исходных величин и знак вопроса; надёжно восстановить расчёт натяжения и скорости по ним невозможно. Полного решения и ответа нет.

\field{Ваш ответ.} Не указан.

\field{Где именно возникает ошибка.} После выписывания данных отсутствует необходимый шаг: второй закон Ньютона для нижней точки с учётом максимального натяжения.

\mistake{} Пропущены уравнение сил и переход к предельной скорости.

\field{Почему это важно.} В нижней точке центр окружности расположен сверху. Натяжение $T$ направлено вверх к центру, сила тяжести $mg$ --- вниз. Поэтому результирующая сила к центру равна $T-mg$. При увеличении скорости требуемое натяжение возрастает; его допустимая граница равна $92$ Н.

\keyidea{} В нижней точке $T-mg=mv^2/R$. Для максимальной скорости подставляют $T=T_{\max}$.

\field{Исправление.} Сначала определите направления обеих сил, затем выразите $v^2$ через предельное натяжение.

\field{Полное правильное решение.}
\begin{align*}
T_{\max}-mg&=m\frac{v_{\max}^2}{R},\\
v_{\max}^2&=\frac{(T_{\max}-mg)R}{m},\\
v_{\max}^2&=\frac{(92-2\cdot10)\cdot4}{2}
=\frac{72\cdot4}{2}=144,\\
v_{\max}&=\boxed{12\ \text{м/с}}.
\end{align*}
\field{Проверка.} При $v=12$ м/с необходимо натяжение $T=mv^2/R+mg=2\cdot144/4+20=92$ Н --- точно допустимый предел.

\result{} Максимально допустимая скорость $12$ м/с.

\field{Задача на закрепление.} Ведёрко массой $3$ кг движется в нижней точке окружности радиусом $2$ м. Нить выдерживает натяжение не более $90$ Н. Найдите максимально допустимую скорость при $g=10\,\text{м/с}^2$.
\clearpage

\taskhead{10}
\field{Условие.} Ведёрко массой $1{,}5$ кг вращается на нити по окружности радиусом $3$ м. Скорость и в верхней, и в нижней точке равна $6$ м/с. Найти оба натяжения и определить, порвётся ли нить, если предельное натяжение $30$ Н.

\field{Ваше решение.} Под номером 10 на фотографии присутствует только знак вопроса. Уравнения для верхней и нижней точек отсутствуют.

\field{Ваш ответ.} Не указан.

\field{Где именно возникает ошибка.} Не выполнено обязательное сравнение двух радиальных балансов сил, а значит невозможно проверить прочность нити.

\mistake{} Отсутствуют расчёты натяжений и сравнение с допустимыми $30$ Н.

\field{Почему это важно.} Центростремительное ускорение всегда направлено к центру окружности. В верхней точке обе силы --- $T_{\text{в}}$ и $mg$ --- направлены к центру. В нижней точке натяжение направлено к центру, а тяжесть --- от центра. Поэтому при одинаковой скорости нижнее натяжение больше верхнего на $2mg$.

\keyidea{} $T_{\text{в}}+mg=mv^2/R$, а $T_{\text{н}}-mg=mv^2/R$.

\field{Исправление.} Составьте два уравнения с осями к центру и сравните оба результата с пределом прочности нити.

\field{Полное правильное решение.} Сначала найдём общую центростремительную силу:
\[
F_{\text{ц}}=m\frac{v^2}{R}=1{,}5\cdot\frac{6^2}{3}=18\ \text{Н},\qquad mg=15\ \text{Н}.
\]
Верхняя точка (ось вниз):
\[
T_{\text{в}}+mg=F_{\text{ц}},\quad
T_{\text{в}}=18-15=\boxed{3\ \text{Н}}.
\]
Нижняя точка (ось вверх):
\[
T_{\text{н}}-mg=F_{\text{ц}},\quad
T_{\text{н}}=18+15=\boxed{33\ \text{Н}}.
\]
\field{Проверка.} Разность $33-3=30$ Н равна $2mg=30$ Н. Верхнее натяжение меньше предела: $3<30$ Н. Нижнее превышает его: $33>30$ Н.

\result{} $T_{\text{в}}=3$ Н, $T_{\text{н}}=33$ Н; нить порвётся в нижней точке при достижении указанной скорости.

\field{Задача на закрепление.} Ведёрко массой $2$ кг вращается на нити длиной $4$ м. В верхней и нижней точках скорость равна $8$ м/с. Найдите оба натяжения и установите, выдержит ли нить предельное натяжение $50$ Н. Примите $g=10\,\text{м/с}^2$.

\end{document}